\documentclass[11pt,a4paper]{article}
\usepackage[margin=2.5cm]{geometry}
\usepackage{amsmath,amssymb,bm}
\usepackage{graphicx}
\usepackage{booktabs}
\usepackage{hyperref}
\usepackage{xcolor}
\usepackage{listings}
\usepackage{algorithm}
\usepackage{algpseudocode}
\usepackage{titlesec}
\usepackage{enumitem}
\usepackage{fancyhdr}
\setlength{\headheight}{14pt}
\usepackage{caption}
\usepackage{subcaption}
\usepackage{tikz}
\usetikzlibrary{arrows.meta,positioning,shapes.geometric,calc}

\definecolor{codeblue}{RGB}{37,99,235}
\definecolor{codegray}{RGB}{107,114,128}
\definecolor{codebg}{RGB}{248,250,252}
\definecolor{accent}{RGB}{37,99,235}

\lstdefinestyle{python}{
  backgroundcolor=\color{codebg},
  basicstyle=\ttfamily\small,
  keywordstyle=\color{codeblue}\bfseries,
  commentstyle=\color{codegray}\itshape,
  stringstyle=\color{green!50!black},
  numbers=left, numberstyle=\tiny\color{codegray},
  numbersep=8pt, frame=single, rulecolor=\color{gray!30},
  breaklines=true, captionpos=b, language=Python,
  showstringspaces=false
}

\pagestyle{fancy}
\fancyhf{}
\lhead{\textcolor{gray}{\small NeuFlow v3 --- Implicit Neural Flow Decoder}}
\rhead{\textcolor{gray}{\small April 27, 2026}}
\cfoot{\thepage}

\hypersetup{colorlinks=true,linkcolor=accent,citecolor=accent,urlcolor=accent}

\title{%
  \vspace{-1cm}
  \textbf{NeuFlow v3: Adapting an Implicit Neural Decoder}\\
  \textbf{from InfiniDepth for Real-Time Optical Flow}\\[0.4em]
  \large Progress Report for PhD Advisor Meeting
}
\author{Shri Raghavendran \\ \small \texttt{github.com/shrirag10/Neuflowv3}}
\date{April 27, 2026}

\begin{document}
\maketitle\thispagestyle{fancy}

%---------------------------------------------------------------
\section{Objective}
%---------------------------------------------------------------
The goal of this project is to replace NeuFlow~v2's \emph{convex upsampler} with an
\emph{implicit neural decoder} inspired by InfiniDepth~\cite{yu2026infinidepth},
improving sub-pixel flow accuracy while keeping the backbone frozen and inference fast.
The key hypothesis: the continuous, coordinate-conditioned MLP can learn finer residuals
than a fixed convex combination of coarse flow values.

%---------------------------------------------------------------
\section{Paper 1: NeuFlow v2 --- Efficient Real-Time Optical Flow}
%---------------------------------------------------------------

\subsection{Core Contribution}
NeuFlow~v2~\cite{neuflowv2} achieves real-time optical flow by replacing the expensive
iterative refinement of RAFT~\cite{raft} with a single-pass architecture:
\begin{enumerate}[leftmargin=*]
  \item \textbf{Multi-scale CNN encoder}: extracts features at $1/8$ and $1/16$ resolution.
  \item \textbf{Cross-attention global correlation} at $1/16$: computes all-pairs matching cheaply using attention instead of 4D cost volumes.
  \item \textbf{GRU-based local refinement} at $1/8$: 8 recurrent iterations refine the flow using local correlation and context features.
  \item \textbf{Convex upsampler}: produces the final $H \times W$ flow as a learned weighted combination of $3\times3$ neighbours from the $1/8$ coarse flow.
\end{enumerate}

\subsection{Convex Upsampler Limitation}
The convex upsampler predicts pixel flow as:
\[
  \mathbf{u}(x,y) = \sum_{i,j \in \mathcal{N}} w_{ij}(x,y)\,\hat{\mathbf{u}}_{ij}
  \quad \text{s.t.} \quad \sum_{i,j} w_{ij}=1,\; w_{ij}\ge 0
\]
where $\hat{\mathbf{u}}$ is the coarse $1/8$ flow and $w$ are learned soft-max weights.
This is a \textbf{fixed discrete mixture} --- it cannot predict flow outside the convex
hull of $3\times3$ neighbours, and it blurs boundaries because nearby cells always mix.

\subsection{Key Architecture Numbers}
\begin{center}
\begin{tabular}{lc}
\toprule
\textbf{Component} & \textbf{Params} \\
\midrule
CNN backbone & 5.1M \\
Cross-attention module & 1.4M \\
GRU refinement & 0.6M \\
Convex upsampler & 0.2M \\
\midrule
Total & 7.3M \\
\bottomrule
\end{tabular}
\end{center}

%---------------------------------------------------------------
\section{Paper 2: InfiniDepth --- Arbitrary-Resolution Depth via Neural Implicit Fields}
%---------------------------------------------------------------

\subsection{Core Contribution}
InfiniDepth~\cite{yu2026infinidepth} represents depth as a \textbf{neural implicit field}:
for any continuous 2D coordinate $(x,y)$, it queries the depth from feature maps instead
of reading a fixed grid. This enables:
\begin{itemize}[leftmargin=*]
  \item Arbitrary output resolution (4K, 8K, 16K from a standard ViT).
  \item Fine-grained boundary prediction (no blurring from convex upsampling).
  \item Differentiable surface normals via autograd.
\end{itemize}

\subsection{Method: Multi-Scale Local Implicit Decoder}
Given input image $I$, a DINOv2-ViT-Large encoder extracts features from layers 4, 11,
and 23, producing a pyramid $\{f^k\}_{k=1}^{L}$ with $f^k \in \mathbb{R}^{h_k \times w_k \times C_k}$
where $(C_1, C_2, C_3) = (256, 512, 1024)$.

\paragraph{Feature Query.}
For any query $(x,y)$, bilinear interpolation at each scale yields local descriptors
$f^k_{(x,y)} \in \mathbb{R}^{C_k}$.

\paragraph{Hierarchical Fusion (Eq.\,3).}
Let $h^1 := f^1_{(x,y)}$ (finest / shallowest scale). For $k = 1, \ldots, L-1$:
\begin{equation}
  h^{k+1} = \mathrm{FFN}^k\!\left(f^{k+1}_{(x,y)} + g^k \odot \mathrm{Linear}(h^k)\right)
  \label{eq:fusion}
\end{equation}
where $g^k \in (0,1)^{C_{k+1}}$ is a \textbf{learnable channel-wise gate} and $\mathrm{FFN}^k$
is a two-layer feed-forward network. Fusion proceeds \emph{shallow $\to$ deep},
integrating fine detail first, then global semantics.

\paragraph{Depth Prediction.}
\[
  d_I(x,y) = \mathrm{MLP}(h^L)
\]

\paragraph{Training (Eq.\,7).}
Supervision on $N$ random coordinate-depth pairs per image:
\[
  \mathcal{L} = \frac{1}{N}\sum_{i=1}^{N}|d_i - \hat{d}_i|
\]
Training: AdamW, lr $=10^{-5}$, 800K steps, 8 GPUs, batch 4.

\subsection{Why Implicit is Better}
\begin{center}
\begin{tabular}{lcc}
\toprule
\textbf{Property} & \textbf{Grid Decoder} & \textbf{Implicit Decoder} \\
\midrule
Output resolution & Fixed (training size) & Arbitrary \\
Boundary handling & Blurred (convex mix) & Sharp (per-point MLP) \\
Sub-pixel accuracy & Limited by stride & Continuous \\
Parameter count & $O(H \cdot W)$ & Fixed (MLP only) \\
\bottomrule
\end{tabular}
\end{center}

%---------------------------------------------------------------
\section{Adaptation: NeuFlow v3 Implicit Decoder}
%---------------------------------------------------------------

\subsection{Conceptual Bridge}
\begin{center}
\begin{tabular}{lll}
\toprule
\textbf{Concept} & \textbf{InfiniDepth} & \textbf{NeuFlow v3} \\
\midrule
Encoder & DINOv2 ViT-L (24 layers) & NeuFlow CNN backbone \\
Feature scales & Layers 4/11/23 (256/512/1024d) & ctx\_s8/feat\_s8/feat\_s16 (64/128/128d) \\
Scale factor & 4$\times$/2$\times$/1$\times$ upsample & 1/8, 1/8, 1/16 native \\
Fusion (Eq.\,\ref{eq:fusion}) & $L=3$ scales, $L-1=2$ steps & Same, $L=3$ \\
Output signal & Scalar depth $d \in \mathbb{R}$ & Flow vector $\Delta\mathbf{u} \in \mathbb{R}^2$ \\
Extra input & None (single image) & Warped img1 features (cross-frame) \\
Loss & Sparse L1 on depth & Sparse L1 on flow residual \\
\bottomrule
\end{tabular}
\end{center}

\subsection{Key Differences from InfiniDepth}
\begin{enumerate}[leftmargin=*]
  \item \textbf{Two-frame correspondence}: We sample img1 features at the warped location
        $(x + u_\text{coarse},\; y + v_\text{coarse})$ and concatenate to the MLP.
        This gives the decoder explicit information about \emph{where img0 moved to}.
  \item \textbf{Residual prediction}: Output is $\Delta\mathbf{u}$, added to bilinearly
        upsampled coarse flow --- not absolute depth.
  \item \textbf{Zero-init}: Output layer weights $\mathbf{W}_\text{out} = \mathbf{0}$,
        so at step 0, $\Delta\mathbf{u} = 0$ and the model behaves as vanilla NeuFlow.
  \item \textbf{Feature dimensions}: CNN backbone caps at 128d vs ViT's 1024d.
        We do not project to higher dims --- no new information would be added.
\end{enumerate}

\subsection{Architecture Diagram}
\begin{center}
\begin{tabular}{ccccc}
\textbf{Scale 1} & & \textbf{Scale 2} & & \textbf{Scale 3} \\
ctx\_s8 (64d) & & feat\_s8 (128d) & & feat\_s16 (128d) \\
\rotatebox{90}{BilinSample$\downarrow$} & & \rotatebox{90}{BilinSample$\downarrow$} & & \rotatebox{90}{BilinSample$\downarrow$} \\
$f_\mathrm{ctx}$ & $\xrightarrow{\sigma(g_1)\odot\mathrm{Lin}}$ & $h^2=\mathrm{FFN1}(f_8 + \cdot)$ & $\xrightarrow{\sigma(g_2)\odot\mathrm{Lin}}$ & $h^3=\mathrm{FFN2}(f_{16} + \cdot)$ \\
\end{tabular}

\bigskip
$\underbrace{[h^3 \mid f_1^\mathrm{warp} \mid \hat{\mathbf{q}} \mid \mathbf{u}_c/WH]}_{260\text{-dim}} \xrightarrow{\mathrm{MLP}_{256\to128\to64\to2}} \Delta\mathbf{u} \xrightarrow{+\,\mathbf{u}_c} \mathbf{u}(x,y)$
\end{center}

\subsection{Full Forward Pass Pseudocode}
\begin{algorithm}[H]
\caption{ImplicitFlowDecoder.forward}
\begin{algorithmic}[1]
\Require feat\_s8 $\in\mathbb{R}^{B\times128\times H/8\times W/8}$, feat1\_s8, feat\_s16 $\in\mathbb{R}^{B\times128\times H/16\times W/16}$,
\Statex \quad\quad ctx\_s8 $\in\mathbb{R}^{B\times64\times H/8\times W/8}$, coarse\_flow $\in\mathbb{R}^{B\times2\times H/8\times W/8}$
\Statex \quad\quad query\_coords $\in\mathbb{R}^{B\times N\times2}$ or \textbf{None} (dense)
\Ensure flow $\in\mathbb{R}^{B\times2\times H\times W}$ or $\mathbb{R}^{B\times N\times2}$

\If{query\_coords is None}
  \State Build dense grid: $\mathbf{q}_{xy} \gets \text{meshgrid}(0\ldots W{-}1,\; 0\ldots H{-}1)$
\EndIf

\State Normalize: $\hat{q}_x \gets \frac{2(q_x+0.5)}{W}-1$,\;\; $\hat{q}_y \gets \frac{2(q_y+0.5)}{H}-1$
\Comment{align\_corners=False}

\State $f_{16} \gets \text{BilinSample}(\texttt{feat\_s16},\;\hat{\mathbf{q}})$
\Comment{$\in\mathbb{R}^{B\times N\times128}$}
\State $f_8 \gets \text{BilinSample}(\texttt{feat\_s8},\;\hat{\mathbf{q}})$
\State $f_\text{ctx} \gets \text{BilinSample}(\texttt{ctx\_s8},\;\hat{\mathbf{q}})$
\Comment{$\in\mathbb{R}^{B\times N\times64}$}

\State \textbf{// 3-scale hierarchical fusion (InfiniDepth Eq.~\ref{eq:fusion}, $L=3$)}
\State $g_1 \gets \sigma(\texttt{gate1})$;\quad $h^2 \gets \texttt{FFN1}\!\left(f_8 + g_1 \odot \texttt{Linear}_{64\to128}(f_\text{ctx})\right)$
\State $g_2 \gets \sigma(\texttt{gate2})$;\quad $h^3 \gets \texttt{FFN2}\!\left(f_{16} + g_2 \odot \texttt{Linear}_{128\to128}(h^2)\right)$

\State $\mathbf{u}_c \gets \text{BilinSample}(\texttt{coarse\_flow},\;\hat{\mathbf{q}})$
\Comment{coarse flow at query pts}
\State $\mathbf{u}_c^{\text{px}} \gets (\mathbf{u}_c^x \cdot W/W_8,\;\; \mathbf{u}_c^y \cdot H/H_8)$

\State \textbf{// Sample img1 features at WARPED position}
\State $\hat{q}^\text{warp}_x \gets \hat{q}_x + 2\mathbf{u}_c^x/W$;\quad $\hat{q}^\text{warp}_y \gets \hat{q}_y + 2\mathbf{u}_c^y/H$
\State $f_1^\text{warp} \gets \text{BilinSample}(\texttt{feat1\_s8},\;\hat{\mathbf{q}}^\text{warp})$

\State $\mathbf{m} \gets [h^3 \;|\; f_1^\text{warp} \;|\; \hat{\mathbf{q}} \;|\; \mathbf{u}_c/WH]$
\Comment{260-dim MLP input}
\State $\Delta\mathbf{u} \gets \texttt{MLP}(\mathbf{m}) \cdot [W, H]$
\Comment{scale back to pixels}

\State $\mathbf{u}(x,y) \gets \mathbf{u}_c^{\text{px}} + \Delta\mathbf{u}$

\If{dense query}
  \State Reshape to $\mathbb{R}^{B\times2\times H\times W}$
\EndIf
\Return $\mathbf{u}$
\end{algorithmic}
\end{algorithm}

\subsection{Training Algorithm}
\begin{algorithm}[H]
\caption{NeuFlow v3 Training Loop}
\begin{algorithmic}[1]
\Require Pretrained NeuFlow checkpoint, $N_\text{steps}=10000$, $\text{lr}=2\times10^{-4}$, $N_\text{pts}=4096$
\State Load backbone $\theta_\text{bb}$; \textbf{freeze} all backbone parameters
\State Initialise decoder $\theta_\text{dec}$; set output layer $\mathbf{W}_\text{out} \gets \mathbf{0}$
\Comment{zero-init: $\Delta\mathbf{u}_0=0$}
\State $\text{optimizer} \gets \text{AdamW}(\theta_\text{dec},\; \text{lr}=2\times10^{-4},\; \text{wd}=10^{-4})$
\For{$t = 1 \ldots N_\text{steps}$}
  \State $(I_0, I_1, \mathbf{u}^*) \sim \mathcal{D}$ \Comment{sample training pair + GT flow}
  \State $\hat{\mathbf{u}}_\text{coarse},\; F \gets \text{Backbone}(I_0, I_1)$
  \Comment{coarse flow + feature maps (frozen)}
  \State \textbf{// Sparse query sampling:} 50\% boundary + 50\% uniform
  \State $\mathbf{q} \gets \text{SampleQueries}(\mathbf{u}^*,\; N_\text{pts})$
  \Comment{$\mathbf{q} \in\mathbb{R}^{N_\text{pts}\times2}$, flow-gradient weighted}
  \State $\hat{\mathbf{u}}(\mathbf{q}) \gets \text{ImplicitDecoder}(F,\; \hat{\mathbf{u}}_\text{coarse},\; \mathbf{q})$
  \State $\mathbf{u}^*(\mathbf{q}) \gets \text{BilinSample}(\mathbf{u}^*,\; \mathbf{q})$
  \Comment{GT flow at query points}
  \State $\mathcal{L} \gets \frac{1}{N_\text{pts}}\|\hat{\mathbf{u}}(\mathbf{q}) - \mathbf{u}^*(\mathbf{q})\|_1$
  \State Backprop; clip gradients ($\|\nabla\|_2 \le 1.0$); step optimizer
  \If{$t \mod 5000 = 0$}
    \State Save checkpoint; log EPE on validation batch
  \EndIf
\EndFor
\end{algorithmic}
\end{algorithm}

%---------------------------------------------------------------
\section{Experiments}
%---------------------------------------------------------------

\subsection{Setup}
\begin{center}
\begin{tabular}{ll}
\toprule
Dataset & VKITTI2 (2,121 train pairs, 1,174 val pairs) \\
Validation scenes & Scene18 + Scene20 \\
Batch size & 4 \\
Sparse points & 4,096 per image \\
Hardware & RTX~4060 Laptop GPU (8\,GB) \\
Training time & $\approx$27 min (10,000 steps) \\
\bottomrule
\end{tabular}
\end{center}

\subsection{Main Results}
\begin{center}
\begin{tabular}{lcccc}
\toprule
\textbf{Model} & \textbf{Mean EPE$\downarrow$} & \textbf{Median EPE$\downarrow$} & \textbf{1px Acc$\uparrow$} & \textbf{3px Acc$\uparrow$} \\
\midrule
NeuFlow v2 (convex upsampler) & 2.23\,px & 1.08\,px & 46.2\% & 78.4\% \\
\textbf{NeuFlow v3 (implicit, 10K)} & \textbf{3.15\,px} & \textbf{1.87\,px} & \textbf{7.2\%} & \textbf{71.1\%} \\
\bottomrule
\end{tabular}
\end{center}

\subsection{Checkpoint Sweep}
\begin{center}
\begin{tabular}{cc}
\toprule
\textbf{Checkpoint} & \textbf{Mean EPE} \\
\midrule
step\_005000 & 3.29\,px \\
\textbf{step\_010000} & \textbf{3.15\,px}~$\leftarrow$~best \\
step\_015000 & 3.47\,px \\
step\_020000 & 3.16\,px \\
step\_025000 & 3.56\,px \\
step\_030000 & 3.69\,px \\
\bottomrule
\end{tabular}
\end{center}

%---------------------------------------------------------------
\section{Analysis and Findings}
%---------------------------------------------------------------

\subsection{Finding 1: Overfitting is the Dominant Bottleneck}
At step 10K, the decoder has seen 2,121 training pairs for $\approx 19$ epochs.
Performance monotonically degrades after 10K:
\[
  \text{steps}/\text{epoch} = \frac{2121}{4} \approx 530 \implies \text{step } 10\text{K} \approx 18.9\text{ epochs}
\]
This is a classic training-set size problem, not an architecture problem.
RAFT, GMFlow, and NeuFlow all train on $>22$K pairs (FlyingChairs) before
fine-tuning --- our training set is $10\times$ too small.

\subsection{Finding 2: End-to-End Training Fails at This Budget}
We attempted joint backbone+decoder training following InfiniDepth's protocol
(backbone lr $=10^{-5}$, decoder lr $=2\times10^{-4}$). The training EPE oscillated
violently between 1.2\,px and 103\,px within 1,000 steps.

\textbf{Root cause:} InfiniDepth requires 800K steps (8 GPUs) for the backbone
and decoder to co-converge. With 30K steps on 1 GPU, the backbone features
shift at every step while the decoder tries to learn from them ---
a \emph{non-stationary learning target} that causes divergence.
Frozen backbone gives stable features and stable convergence.

\subsection{Finding 3: 3-Scale Fusion Helps Marginally}
A 2-scale decoder (without ctx\_s8) trained for 20K steps: EPE\,$=3.20$\,px. \\
The 3-scale decoder trained for 10K steps: EPE\,$=3.15$\,px. \\
Consistent with InfiniDepth's finding that hierarchical fusion benefits most
when the decoder is trained to convergence on sufficient data.

%---------------------------------------------------------------
\section{Next Steps}
%---------------------------------------------------------------

\begin{enumerate}[leftmargin=*]
  \item \textbf{FlyingChairs + FlyingThings3D} (highest priority): adds $22{,}872 + 21{,}818 = 44{,}690$ training pairs. Expected to remove the overfitting bottleneck and allow the decoder to converge, targeting EPE $< 2.23$\,px.

  \item \textbf{Conservative gate initialisation}: set \texttt{gate} $= -2$ so $\sigma(-2)\approx0.12$ initially, allowing stable warm-up before the context features contribute.

  \item \textbf{Phase-2 end-to-end fine-tuning}: once decoder converges with frozen backbone on the full curriculum, fine-tune the full model at lr $=5\times10^{-6}$ for 20K steps. This is the standard 2-phase transfer protocol and should further reduce EPE.
\end{enumerate}

%---------------------------------------------------------------
\section{Conclusion}
%---------------------------------------------------------------
We have successfully implemented an InfiniDepth-style implicit decoder for NeuFlow v2.
The prototype achieves 3.15\,px EPE with only 265K new parameters (3.6\% of model)
and 27 minutes of training. The architecture and training protocol are sound ---
the primary limitation is dataset size, a known and solvable constraint.
The code is available at: \url{https://github.com/shrirag10/Neuflowv3}.

\begin{thebibliography}{9}
\bibitem{yu2026infinidepth}
Hao Yu et al., \emph{InfiniDepth: Arbitrary-Resolution and Fine-Grained Depth Estimation with Neural Implicit Fields}, arXiv:2601.03252, 2026.

\bibitem{neuflowv2}
Shengyu Zhao et al., \emph{NeuFlow v2: High-Efficiency Optical Flow Estimation on Edge Devices}, arXiv:2408.10161, 2024.

\bibitem{raft}
Zachary Teed and Jia Deng, \emph{RAFT: Recurrent All-Pairs Field Transforms for Optical Flow}, ECCV 2020.
\end{thebibliography}

\end{document}
